\onecolumn

\section{Benchmark Presets}
\label{appendix:presets}

\begin{table}[h]
\centering
\setlength{\tabcolsep}{4pt}
\renewcommand{\arraystretch}{1.05}
\resizebox{0.65\textwidth}{!}{%
\begin{tabular}{@{}lccccc@{}}
\hline
\textbf{Preset} & \textbf{n\_layer} & \textbf{n\_embd} & \textbf{block\_size} & \textbf{n\_head} & \textbf{steps} \\
\hline

\hline
S  & 1 & 64  & 128 & 4  & 2000  \\
M  & 2 & 128 & 128 & 8  & 5000  \\
L  & 4 & 256 & 256 & 8  & 10000 \\
VL & 6 & 384 & 512 & 12 & 20000 \\
XL & 8 & 512 & 512 & 16 & 30000 \\

\hline
MS5          & 1 & 64  & 512 & 4  & 5000 \\
MM5          & 2 & 128 & 512 & 8  & 5000 \\
ML5          & 4 & 256 & 512 & 8  & 5000 \\
MVL5         & 6 & 384 & 512 & 12 & 5000 \\
MXL5         & 8 & 512 & 512 & 16 & 5000 \\

\hline
\end{tabular}%
}
\caption{Model Configurations Used in Experiments}
\label{tab:benchmark_presets}
\end{table}

\section{Serial Implementation Prompt}
\label{appendix:serial-prompt}
The initial serial C++ baseline discussed in Section~2 was generated from the prompt ``Convert the provided microGPT-style Python implementation into a clean serial C++ program,'' with the following requirements:
\begin{itemize}[noitemsep, topsep=0pt]
    \item Keep the model architecture equivalent to the Python version, including token embeddings, positional embeddings, stacked transformer blocks, RMSNorm, causal self-attention, MLP feedforward layers, residual connections, and final logits with cross-entropy loss.
    \item Implement only the forward pass.
    \item Do not use CUDA, OpenMP, Eigen, BLAS, or any external numerical library.
    \item Use simple C++ data structures such as \texttt{std::vector} over \texttt{double} values and explicit loops.
    \item Process one input sequence at a time.
    \item Preserve the same tensor shapes and operation ordering as the Python version.
    \item Include clear helper functions for embedding lookup, RMSNorm, linear projection, attention score computation, softmax, MLP projection, ReLU, residual addition, and logits/loss computation.
    \item Prioritize readability and correctness over optimization.
    \item Use deterministic parameter loading from fixture files so the C++ implementation can be checked against the Python baseline.
    \item Organize the code so later versions can reuse the same structure for batching and CUDA parallelization.
\end{itemize}

\section{Optimized Profiling Results}
\label{appendix:profile}
\begin{table}[h]
\centering
\setlength{\tabcolsep}{4pt}
\renewcommand{\arraystretch}{1.05}
\resizebox{0.65\textwidth}{!}{%
\begin{tabular}{@{}lccccc@{}}
\hline
\textbf{Kernel} & \textbf{Time (\%)} & \textbf{Total Time} & \textbf{Calls} & \textbf{Avg} & \textbf{Max} \\
\hline
linear\_kernel & 85.88\% & 2.40485 s & 1372 & 1.7528 ms & 6.9377 ms \\
rmsnorm\_kernel & 4.52\% & 126.59 ms & 476 & 265.95 $\mu$s & 424.22 $\mu$s \\
self\_attn\_kernel & 4.45\% & 124.61 ms & 224 & 556.29 $\mu$s & 1.1746 ms \\
cross\_entropy\_loss\_kernel & 3.73\% & 104.36 ms & 28 & 3.7271 ms & 6.7604 ms \\
relu\_kernel & 0.58\% & 16.380 ms & 224 & 73.123 $\mu$s & 133.66 $\mu$s \\
add\_vec\_kernel & 0.42\% & 11.762 ms & 448 & 26.254 $\mu$s & 50.016 $\mu$s \\
embedding\_lookup\_kernel & 0.02\% & 481.57 $\mu$s & 28 & 17.198 $\mu$s & 31.776 $\mu$s \\
\hline
\end{tabular}%
}
\label{tab:profile_optimization}
\caption{CUDA Kernel Profiling Results for the Optimized Parallel Implementation}
\end{table}
